\scanpage{86}
\mspage{73}

la stabilité de $W$ par $\unsure{\varinjlim|\underline{\lim}}$ filtrantes (donc aussi par facteurs directs) assure alors qu'on a également
\[
C\subset\mathrm{Cof}_W,\tag*{(4.7)}
\]
et par suite
\[
\left\{\begin{aligned} C\cap W&\subset\mathrm{Cof}_W\cap W=W^{\mathrm{univ}}\\ \text{d'où}\quad \mathrm{TC}&\overset{\mathrm{déf}}{=}C\cap W=C\cap W^{\mathrm{univ}},
\end{aligned}\right.\tag*{(4.8)}
\]
ce qui assure la stabilité de $C\cap W$ par cochg. de base. La stabilité \inserted{de $\mathrm{TC}$} par comp. transfinie résulte de celle de $C$, et de la stabilité de \deleted{\unsure{$C$}}$W$ par comp. transfinie.

4.3. Je reviens aux conditions \unsure{primitives|générales} \textcircled{a} \unsure{à|et} \textcircled{\unsure{f|g}} de 4.1, sans supposer néc. $W$ stable par $\varinjlim$ filtrantes. Je veux à présent \unsure{étendre|étudier} la structure à cofibrations \unsure{aux|sur} $M(I)$, $I$ quelc. de $\mathrm{Cat}$. Nous \unsure{n'avons|n'aurons} pas pour ceci \inserted{(je crois)} à utiliser les développements des \S\S\,2,3, mais nous ferons le lien avec ceux-ci par la suite. \tikz[x=1ex,y=1ex]{\draw (0,0)--(.9,0)--(1.26,1.35); \draw (.8,.08)--(1.15,1.4); \draw (1.35,1.43)--(2.2,1.6);}

Désignons par $I_0$ la catégorie discrète ayant mêmes objets que $I$, d'où un foncteur canonique
\[
\unsure{u|v}=u_I:I_0\longrightarrow I.\tag*{(4.9)}
\]
L'idée est de définir la structure de $M(I)$ via la structure évidente de $M(\unsure{I_0})\simeq M\deleted{\unsure{\ell}}^{I_0}$, en jouant sur le couple de foncteurs adjoints
\diagram{\begin{tikzcd}[ampersand replacement=\&,column sep=huge] M(I_0) \arrow[r,shift left,"\varphi=u_!"] \& M(I) \arrow[l,shift left,"\psi=u^*"]
\end{tikzcd}\,.}
% archive sheet number: 85
